% TOP_PROBLEM: {"id": 7200088, "title": "The Thomson problem", "category_id": 6, "category": {"id": 6, "name": "geometry", "display_name": "Geometry", "description": "Euclidean and non-Euclidean geometry, geometric structures.", "slug": "geometry", "order_index": 6, "created_at": "2026-07-31T15:26:25.671Z"}, "status": "partially_solved"}
% FIRST_POSTED: 2026-09-27
% !TeX program = lualatex
% ATTEMPT_STATUS: unresolved
% Model: GPT-6 Astra Ultra; continuation dated 27 September 2026.
\documentclass[11pt]{article}
\usepackage[margin=25mm]{geometry}
\usepackage{fontspec}
\setmainfont{Latin Modern Roman}
\usepackage{amsmath,amssymb,mathtools}
\usepackage{unicode-math}
\setmathfont{Latin Modern Math}
\usepackage{xurl,hyperref}
\hypersetup{colorlinks=true,urlcolor=blue}
\setcounter{secnumdepth}{0}
\setlength{\parindent}{0pt}
\setlength{\parskip}{5pt}
\setlength{\emergencystretch}{3em}
\begin{document}
\section*{\texorpdfstring{The Thomson problem}{The Thomson problem}}\label{7200088}
\subsection{Definitions and mathematical statement}
For $N\ge1$, choose pairwise distinct points $x_1,\ldots,x_N$ on the unit sphere $S^2\subset{\mathbb{R}}^3$. In units with equal unit charges and Coulomb constant one, the interaction energy is
\[
 E_N(x_1,\ldots,x_N)=\sum_{1\le i<j\le N}\frac1{\|x_i-x_j\|_2}.
\]
The Thomson problem asks to determine minimizing configurations for every $N$, up to rigid motions and permutation of the charges. Distances are Euclidean chord distances, not spherical arc lengths. Coincident charges have infinite energy. Existence of a minimizer or a numerical configuration with low energy does not give its rigorous global characterization; degeneracy of minimizers is allowed, and no uniqueness claim is built into the question.
\par\smallskip\textit{Formulation and concept references:} \href{https://web.cels.anl.gov/~zippy/publications/thomson/thomsonPRL.html}{[S1]}.

\subsection{Short English statement}
For every number of equal charges, which arrangements on the unit sphere minimize their total Coulomb repulsion?
\subsection{Research attempt}
\paragraph{Scope and method.}
The inherited source [S1] concerns the Coulomb problem and the
difficulty of identifying global minimizing configurations. A low
numerical value is only an upper bound on the minimum. The present
attempt supplies exact certificates for $N\leq4$ and $N=6$, a
first- and second-variation test, and a global leading-order estimate.
These do not classify the minimizing arrangements for arbitrary $N$.
All distances below are chord distances, as in the selected statement.

\paragraph{Existence without an accidental uniqueness claim.}
Extend the energy to $(S^2)^N$ by the value $+\infty$ at collisions.
Each pair potential is lower semicontinuous, and so is their finite
sum. The product of spheres is compact, hence the extended energy
attains its minimum. There are distinct configurations with finite
energy, so a minimizer contains no collision. For $N=1$ the minimum
is zero and every point is equivalent under rotation.

If a finite-energy configuration has energy at most $U$, every term
in its positive energy sum is at most $U$. Thus
\[
 |x_i-x_j|\geq U^{-1}\quad(i\ne j).
 \tag{391.1}
\]
This elementary separation bound is sufficient to put a bounded-energy
search in a compact collision-free region. It is generally too weak
to identify an optimal configuration, but prevents a numerical search
from being justified using a compactness statement that ignores
collisions.

\paragraph{A centroid bound sharp for the small simplexes.}
Let $M=N(N-1)/2$ and $s=\sum_i x_i$. Direct expansion yields
\[
 \sum_{i<j}|x_i-x_j|^2=N\sum_i|x_i|^2-|s|^2=N^2-|s|^2\leq N^2.
 \tag{391.2}
\]
The function $r\mapsto r^{-1/2}$ is strictly convex and decreasing
on $(0,\infty)$. Jensen's inequality and (391.2) therefore give
\[
 E_N\geq M\left(\frac1M\sum_{i<j}|x_i-x_j|^2\right)^{-1/2}
     \geq\frac{M^{3/2}}N.
 \tag{391.3}
\]
Equality requires both $s=0$ and equal squared distances. For
$N=2,3,4$, these requirements are realized by an antipodal pair,
a regular equilateral triangle on a great circle, and a regular
tetrahedron. Their exact energies are respectively
\[
 \frac12,\qquad\sqrt3,\qquad 3\sqrt{3/2}.
 \tag{391.4}
\]
This also proves the characterization in these cases: the Gram matrix
has diagonal one and off-diagonal $-1/(N-1)$, and equal Gram matrices
are related by an orthogonal map on their linear spans, extending to
an orthogonal map of $\mathbb R^3$.

For $N>4$ the equality Gram matrix has rank $N-1>3$, so equality in
(391.3) is impossible. This is an explicit obstruction to extrapolating
the regular-simplex proof: the desired equidistant point set does not
fit in the ambient space. A strict inequality without a quantitative
gap gives no replacement characterization.

\paragraph{A quadratic certificate proving the octahedron case.}
Write $t=x_i\cdot x_j$ and $f(t)=(2-2t)^{-1/2}$ for $-1\leq t<1$.
Consider the polynomial
\[
 P(t)=a+bt+ct^2,\qquad
 a=\frac1{\sqrt2},\quad b=\frac1{2\sqrt2},\quad
 c=\frac12-\frac1{2\sqrt2}>0.
 \tag{391.5}
\]
It has $P(-1)=f(-1)$, $P(0)=f(0)$ and $P'(0)=f'(0)$.
In fact
\[
 P(t)\leq f(t)\quad(-1\leq t<1),
 \quad\hbox{with equality only at }t=-1,0.
 \tag{391.6}
\]
Here is a direct remainder proof. For a fixed $t$ distinct from the
two interpolation nodes, subtract a suitable multiple
$A(s+1)s^2$ from $f(s)-P(s)$ to make it vanish also at $s=t$.
There are four zeros counting the double zero at zero, so repeated
Rolle's theorem gives $6A=f'''(\xi)>0$ for an intermediate $\xi$.
Therefore $f(t)-P(t)=A(t+1)t^2>0$. At the nodes equality is already
built in. This proof works also for $-1<t<0$; the sign of the
remainder is not guessed from a graph.

Set $Q=\sum_i x_ix_i^{\mathsf T}$. Then $\operatorname{tr}Q=N$,
and
\[
 \sum_{i<j}t_{ij}=\tfrac12(|s|^2-N),\qquad
 \sum_{i<j}t_{ij}^2=\tfrac12(\operatorname{tr}(Q^2)-N).
\]
Since the three eigenvalues of $Q$ are nonnegative, Cauchy--Schwarz
gives $\operatorname{tr}(Q^2)\geq N^2/3$. Combining these identities
with $b,c>0$ proves the general lower bound
\[
 E_N\geq aM-\frac{bN}{2}+
                  \frac c2\left(\frac{N^2}{3}-N\right).
 \tag{391.7}
\]
At $N=6$ the right side is $6\sqrt2+3/2$. The six vertices
$\{\pm e_1,\pm e_2,\pm e_3\}$ attain it: there are three antipodal
pairs of distance two and twelve remaining pairs of distance $\sqrt2$.
Thus the octahedron is a global minimizer with this exact energy.

The equality conditions characterize it. Equality in (391.6) forces
every inner product between distinct vertices to be zero or minus
one. Distinct unoriented lines represented by vertices must therefore
be mutually orthogonal, so there are at most three such lines. Each
line contains at most two unit points. Six points force exactly three
orthogonal lines with both signs present. This is the octahedron up
to rigid motion and relabeling. The argument uses positivity of the
quadratic coefficients to certify a \emph{global} bound; it is
stronger than checking a Hessian at the candidate.

\paragraph{What stationarity and the Hessian actually test.}
For a collision-free critical configuration, tangential projection
of the ordinary gradient must vanish:
\[
 (I-x_ix_i^{\mathsf T})
       \sum_{j\ne i}\frac{x_i-x_j}{|x_i-x_j|^3}=0
       \quad(1\leq i\leq N).
 \tag{391.8}
\]
The minus sign in the gradient is immaterial after setting the
projection to zero. Symmetry often proves (391.8), but it does not
compare energies of different critical configurations.

Let $v_i\perp x_i$ and move each point on its great-circle geodesic
with initial velocity $v_i$. At zero time its acceleration is
$-|v_i|^2x_i$. With $a_{ij}=x_i-x_j$, $b_{ij}=v_i-v_j$ and
$d_{ij}=|a_{ij}|$, differentiation gives the second variation
\[
 \sum_{i<j}\left[
 \frac{3(a_{ij}\cdot b_{ij})^2}{d_{ij}^5}
 -\frac{|b_{ij}|^2-(1-x_i\cdot x_j)(|v_i|^2+|v_j|^2)}{d_{ij}^3}
 \right].
 \tag{391.9}
\]
To verify the curvature correction, differentiate
$|x_i(t)-x_j(t)|^2$ twice: its second derivative is
$2|b_{ij}|^2-2(1-x_i\cdot x_j)(|v_i|^2+|v_j|^2)$.
Applying the chain rule to the inverse square root gives (391.9).
At a local minimum this quadratic form is nonnegative. Global
rotations produce zero directions, which must not be interpreted as
instability. Conversely, a negative direction disproves local
minimality. Nonnegativity alone proves neither uniqueness nor global
minimality, and a floating-point nearly zero eigenvalue is not an
exact certificate of either sign.

\paragraph{A uniform upper bound from random configurations.}
Fix one point $x$ on the sphere and let $Y$ be uniform on $S^2$.
The inner product $t=x\cdot Y$ is uniform on $[-1,1]$. Therefore
\[
 \mathbb E\frac1{|x-Y|}
   =\frac12\int_{-1}^1(2-2t)^{-1/2}\,dt=1.
\]
For $N$ independent uniform points, collisions have probability zero
and $\mathbb E E_N=M$. Hence some configuration has energy at most
$M$, and the minimum satisfies $E_N^{\min}\leq N(N-1)/2$.
Expectation is used only for existence; it gives no explicit
minimizing configuration.

\paragraph{A matching leading-order lower bound by smearing charges.}
We give the positivity detail needed in an electrostatic lower bound.
For finite signed compactly supported measures of finite Newton
energy, the kernel $1/|x-y|$ is positive definite. One verification
uses
\[
 \frac1r=\frac1{\sqrt\pi}\int_0^\infty t^{-1/2}e^{-tr^2}\,dt.
\]
Each Gaussian kernel is positive definite, by its positive Fourier
transform, so its quadratic energy is nonnegative. Truncating the
$t$ integral and then passing to the limit proves the assertion;
finite absolute energy justifies the limit for the signed measures
used here. Equivalently one can first regularize them smoothly.

Let $\sigma_{x,\varepsilon}$ be probability surface measure on the
radius-$\varepsilon$ sphere centered at $x$, and put
$\nu=\sum_i\sigma_{x_i,\varepsilon}$. A direct angular integral
gives the shell potential
\[
 \int\frac{d\sigma_{x,\varepsilon}(y)}{|z-y|}
   =\frac1{\max\{|z-x|,\varepsilon\}}.
 \tag{391.10}
\]
The self-energy of each shell is $1/\varepsilon$. Its potential
is at most the point-charge potential. Averaging that inequality
over a second shell and using (391.10) again shows that the mutual
energy of two shells is at most $1/|x_i-x_j|$, even if the shells
overlap. Consequently
\[
 I(\nu):=\iint\frac{d\nu(x)d\nu(y)}{|x-y|}
      \leq 2E_N+N/\varepsilon.
 \tag{391.11}
\]
All of $\nu$ is supported in the radius-$R$ ball with $R=1+\varepsilon$.
Let $\sigma_R$ be uniform probability measure on its boundary sphere.
Its potential equals $1/R$ throughout that ball. Positivity of the
energy of $\nu-N\sigma_R$ therefore yields
\[
 0\leq I(\nu-N\sigma_R)=I(\nu)-N^2/R.
\]
Together with (391.11), this gives
\[
 E_N\geq\frac{N^2}{2(1+\varepsilon)}-\frac N{2\varepsilon}.
 \tag{391.12}
\]
Choose $\varepsilon=N^{-1/2}$ and use
$(1+\varepsilon)^{-1}\geq1-\varepsilon$. We obtain the entirely
uniform bounds
\[
 \frac{N^2}{2}-N^{3/2}\leq E_N^{\min}\leq\frac{N(N-1)}2.
 \tag{391.13}
\]
In particular $E_N^{\min}/N^2\to1/2$. These bounds leave a large
configuration-sensitive next-order term undetermined. Leading
asymptotics alone cannot discriminate between competing finite
arrangements whose energies differ far below order $N^2$.

\paragraph{Diagnostics, certificates and the first gap.}
The accompanying checks compare floating-point energies for the small
candidates with the derived formulas, numerically check the polynomial
contacts and sampled minorant values, and compare the second-variation
formula with finite differences away from collisions. The analytic inequality
(391.6), rather than a sampled graph, certifies the global lower
bound for $N=6$.

One possible route for another $N$ is a minorizing polynomial with
nonnegative harmonic coefficients and contacts at all inner products
of a proposed configuration. Such a polynomial would need a verified
inequality on the entire interval $[-1,1)$ and equality in every
resulting moment bound. There is no reason that an arbitrary candidate
admits a certificate of this particular form. Failure to find one
does not establish that the candidate is wrong, and local numerical
optimization does not replace one.

The results above settle $N=1,2,3,4,6$ within this attempt and give
uniform leading asymptotics. The first gap toward the catalogue
question is a rigorous global comparison for general configurations
at every remaining $N$. No such classification follows from the
centroid, Hessian or smearing bounds. The all-$N$ problem remains
unresolved by this attempt.
\subsection{Sources}
\begin{itemize}
\item[S1]Altschuler et al., The Thomson Problem, Physical Review Letters author manuscript page (1997).
\url{https://web.cels.anl.gov/~zippy/publications/thomson/thomsonPRL.html}.
\textit{Formulation source.}
\end{itemize}
\end{document}
